
@@M0@@ a une
résolution partielle globale~:
@@M1@@
par des @@M2@@-Modules localement libres de type fini, le noyau @@M3@@ de @@M4@@ est encore un @@M5@@-Module de type fini (I. 2.11). Il est même de présentation finie (en faisant le même raisonnement avec @@M6@@ au lieu de @@M7@@). Enfin, il est plat sur @@M8@@. Il vérifie donc les mêmes hypothèses que @@M9@@, ce qui permet de construire de proche en proche une résolution globale graduée de @@M10@@ par des @@M11@@-Modules gradués localement libres de type fini, d'après la remarque initiale. Pour @@M12@@, @@M13@@ est un @@M14@@-Module projectif sur tout ouvert affine (1.5.2) et par suite est localement libre en tant que @@M15@@-Module gradué, d'après 1.5.1. On peut alors, dans la résolution graduée construite, remplacer @@M16@@ par @@M17@@, et les termes suivants par @@M18@@, d'où la résolution annoncée.

Montrons enfin le point c) de la démonstration.

\textbf{Proposition 1.7.} \uline{Soient} @@M19@@ \uline{un schéma quasi-compact}, @@M20@@ \uline{un} @@M21@@\uline{-Module localement libre de type fini}, @@M22@@, \uline{et} @@M23@@ \uline{un} @@M24@@\uline{-Module gradué}, \uline{localement libre de type fini} (1.2). \uline{Alors il existe une filtration décroissante de} @@M25@@ \uline{par des sous-}@@M26@@\uline{-Modules gradués} @@M27@@, \uline{telle que} :
@@M28@@

Comme @@M29@@ est quasi-compact, on peut trouver des entiers @@M30@@ et @@M31@@ tels qu'au voisinage de chaque point de @@M32@@, les éléments d'une base homogène 
